\chapter{Plantations}
\label{chap:plantations}

\section{Vue d'ensemble}

Une \emph{plantation} (\emph{planting}) représente la mise en culture d'une variété à un endroit donné. Pomone distingue deux grandes catégories :

\begin{itemize}
  \item Les \textbf{plantations annuelles} : un cycle de vie complet sur une saison (semis $\to$ récolte $\to$ retrait). Typique du maraîchage diversifié.
  \item Les \textbf{plantations pluriannuelles} : une plantation qui reste en place plusieurs années (vergers, petits fruits, asperges, rhubarbe, vivaces). Les récoltes sont enregistrées année par année.
\end{itemize}

\section{Créer une plantation}

\begin{todobox}
À détailler : formulaire de création, sélection variété/lieu, dates clés (semis serre, mise en place, retrait), surface, nombre de plants.
\end{todobox}

\section{Vue détail d'une plantation}

\begin{todobox}
À détailler : informations affichées (variété, lieu, surface, calendrier, état), boutons d'action.
\end{todobox}

\section{Récoltes annuelles pour plantations pluriannuelles}

\begin{todobox}
À détailler : tableau des récoltes par année, ajout d'une récolte, comparaison estimé/réalisé.
\end{todobox}

\section{Vue Gantt des plantations}

\begin{todobox}
\textbf{À venir.} Vue chronologique de toutes les plantations de la saison sous forme de barres horizontales (semis $\to$ champ $\to$ récolte), avec ligne du jour, drag-and-drop, et superposition des tâches.
\end{todobox}
